%% ----------------------------------------------------------------------------
%%  §6.4  Funzioni convesse e concave
%% ----------------------------------------------------------------------------
\section{Funzioni convesse e concave}
\label{sec:convessita}

Il segno della derivata prima dice se il grafico sale o scende; non dice però
\emph{come} sale o scende. Le funzioni $x^{2}$ e $\sqrt x$ sono entrambe crescenti in
$[0,+\infty)$, ma la prima cresce sempre più rapidamente e il suo grafico volge la
concavità verso l'alto, la seconda cresce sempre più lentamente e il suo grafico volge
la concavità verso il basso. Per descrivere questa proprietà introduciamo una nuova
definizione.

Si dice che una funzione è \emph{convessa} in un intervallo se, per ogni punto
$x_{0}$ dell'intervallo, il grafico della funzione sta \emph{al di sopra} della retta
tangente nel punto $(x_{0},f(x_{0}))$; si dice \emph{concava} se il grafico sta \emph{al di
sotto} di ogni retta tangente. Usando l'equazione della retta tangente
\eqref{eq:retta-tangente}, possiamo dare la definizione in forma precisa.

\begin{definizione}[Funzioni convesse e concave]\label{def:convessa}
\index{funzione!convessa}\index{funzione!concava}
Sia $f$ una funzione derivabile nell'intervallo $[a,b]$ (con derivate unilaterali negli
estremi). Si dice che $f$ è \emph{convessa} in $[a,b]$ se
\begin{equation}\label{eq:convessa}
  f(x)\ge f(x_{0})+f'(x_{0})(x-x_{0}) \qquad\perogni x,x_{0}\in[a,b];
\end{equation}
si dice che $f$ è \emph{concava} in $[a,b]$ se
\begin{equation}\label{eq:concava}
  f(x)\le f(x_{0})+f'(x_{0})(x-x_{0}) \qquad\perogni x,x_{0}\in[a,b].
\end{equation}
Se le disuguaglianze sono strette per $x\neq x_{0}$, la funzione si dice
\emph{strettamente} convessa o concava. Una funzione è convessa (concava) in un
intervallo qualunque $I$ se lo è in ogni intervallo chiuso e limitato contenuto in $I$.
\end{definizione}

Evidentemente $f$ è concava se e solo se $-f$ è convessa: per questo basta
studiare le funzioni convesse. Nella \cref{fig:convessa-concava} la funzione è convessa
nell'intervallo $[a,c]$ e concava nell'intervallo $[c,b]$: in $c$ la concavità cambia
verso.

\begin{definizione}[Punto di flesso]\label{def:flesso}\index{punto!di flesso}\index{flesso}
Sia $f$ derivabile in un intorno di $x_{0}$. Si dice che $x_{0}$ è un \emph{punto di
flesso} per $f$ se esiste $\delta>0$ tale che $f$ è convessa in $[x_{0}-\delta,x_{0}]$ e
concava in $[x_{0},x_{0}+\delta]$, oppure concava in $[x_{0}-\delta,x_{0}]$ e convessa in
$[x_{0},x_{0}+\delta]$.
\end{definizione}

In un punto di flesso la retta tangente \emph{attraversa} il grafico: da una parte il
grafico sta sopra la tangente, dall'altra sotto. Se la tangente è orizzontale si parla
di \emph{flesso a tangente orizzontale} (come per $x^{3}$ in $0$), altrimenti di flesso a
tangente obliqua.

\begin{figure}[htb]
\centering
\begin{tikzpicture}
\begin{axis}[width=0.9\textwidth, height=5.6cm, xmin=0.4, xmax=5.6, ymin=1.15, ymax=2.85,
  axis lines=left, xtick={0.8,3,5.2}, xticklabels={$a$,$c$,$b$}, ytick=\empty,
  clip=false]
  \addplot[curva, domain=0.8:5.2, samples=120] {2-0.15*(x-3)^3+0.5*(x-3)};
  \addplot[curva2, domain=0.55:2.0] {1.9748-0.958*(x-1.2)};
  \addplot[curva2, domain=2.2:3.8] {2+0.5*(x-3)};
  \addplot[curva2, domain=4.0:5.45] {2.0252-0.958*(x-4.8)};
  \addplot[tratteggio] coordinates {(3,1.15) (3,2)};
  \addplot[only marks, mark=*, mark size=1.6pt, blu] coordinates
    {(1.2,1.9748) (3,2) (4.8,2.0252)};
  \node[font=\scriptsize, grigio] at (axis cs:1.9,2.65) {convessa};
  \node[font=\scriptsize, grigio] at (axis cs:3.6,2.7) {concava};
  \node[font=\scriptsize, rame, anchor=north west] at (axis cs:3.05,1.95) {flesso};
\end{axis}
\end{tikzpicture}
\caption{Una funzione convessa in $[a,c]$, dove il grafico sta sopra le tangenti, e
concava in $[c,b]$, dove sta sotto. Nel punto di flesso $c$ la tangente attraversa il
grafico.}
\label{fig:convessa-concava}
\end{figure}

Come per la monotonia, esiste un criterio che riconduce la convessità al segno di una
derivata, questa volta della derivata seconda.

\needspace{12\baselineskip}
\begin{teorema}[Criterio di convessità]\label{teo:criterio-convessita}
\index{criterio!di convessità}
Sia $f$ una funzione derivabile in $[a,b]$ che ammette derivata seconda in $(a,b)$. Le
seguenti condizioni sono fra loro equivalenti:
\begin{enumerate}[label=\textup{(\alph*)}]
  \item $f$ è convessa in $[a,b]$;
  \item $f'$ è crescente in $[a,b]$;
  \item $f''(x)\ge0$ per ogni $x\in(a,b)$.
\end{enumerate}
Analogamente, $f$ è concava in $[a,b]$ se e solo se $f'$ è decrescente in $[a,b]$, se e
solo se $f''\le0$ in $(a,b)$.
\end{teorema}

\begin{schema}
\Ipotesi{$f$ derivabile in $[a,b]$, derivabile due volte in $(a,b)$.}
\Tesi{(a) $f$ convessa \ $\sse$ \ (b) $f'$ crescente \ $\sse$ \ (c) $f''\ge0$ in $(a,b)$.}
\end{schema}

\begin{strategia}
Dimostriamo le equivalenze (a)$\sse$(b) e (b)$\sse$(c).
\begin{itemize}[nosep]
  \item (a)$\Rightarrow$(b): scriviamo la disuguaglianza di convessità con il punto di
  tangenza in $x_{1}$ e poi in $x_{2}$, e sommiamo: i valori di $f$ si cancellano e resta
  una disuguaglianza tra le derivate.
  \item (b)$\Rightarrow$(a): per il teorema di Lagrange la differenza $f(x)-f(x_{0})$ è
  uguale a $f'(x_{1})(x-x_{0})$, con $x_{1}$ tra $x_{0}$ e $x$; la monotonia di $f'$
  permette di confrontare $f'(x_{1})$ con $f'(x_{0})$.
  \item (b)$\sse$(c): è il criterio di monotonia applicato alla funzione $f'$, con una
  piccola cautela negli estremi $a$ e $b$, dove non sappiamo se $f'$ sia continua.
\end{itemize}
\end{strategia}

\begin{dimostrazione}
\emph{(a)$\Rightarrow$(b).} Siano $x_{1},x_{2}\in[a,b]$ con $x_{1}<x_{2}$. Scrivendo la
\eqref{eq:convessa} con $x_{0}=x_{1}$, $x=x_{2}$, e poi con $x_{0}=x_{2}$, $x=x_{1}$,
otteniamo
\[
  f(x_{2})\ge f(x_{1})+f'(x_{1})(x_{2}-x_{1}), \qquad
  f(x_{1})\ge f(x_{2})+f'(x_{2})(x_{1}-x_{2}) .
\]
Sommando membro a membro e semplificando,
\begin{equation}\label{eq:convessa-somma}
  0\ge f'(x_{1})(x_{2}-x_{1})+f'(x_{2})(x_{1}-x_{2})=-\bigl[f'(x_{2})-f'(x_{1})\bigr](x_{2}-x_{1}),
\end{equation}
cioè $\bigl[f'(x_{2})-f'(x_{1})\bigr](x_{2}-x_{1})\ge0$. Essendo $x_{2}-x_{1}>0$, ne segue
$f'(x_{1})\le f'(x_{2})$: $f'$ è crescente.

\emph{(b)$\Rightarrow$(a).} Siano $x,x_{0}\in[a,b]$ con $x\neq x_{0}$. Per il teorema di
Lagrange (\cref{teo:lagrange}) esiste $x_{1}$ strettamente compreso tra $x_{0}$ e $x$
tale che
\begin{equation}\label{eq:convessa-lagrange}
  f(x)-f(x_{0})=f'(x_{1})(x-x_{0}) .
\end{equation}
Se $x>x_{0}$, allora $x_{1}>x_{0}$ e, per la monotonia di $f'$, $f'(x_{1})\ge f'(x_{0})$;
moltiplicando per $x-x_{0}>0$ si ottiene $f'(x_{1})(x-x_{0})\ge f'(x_{0})(x-x_{0})$. Se
$x<x_{0}$, allora $x_{1}<x_{0}$ e $f'(x_{1})\le f'(x_{0})$; moltiplicando per
$x-x_{0}<0$ il verso si inverte, e si ottiene di nuovo
$f'(x_{1})(x-x_{0})\ge f'(x_{0})(x-x_{0})$. In entrambi i casi, per la
\eqref{eq:convessa-lagrange}, $f(x)-f(x_{0})\ge f'(x_{0})(x-x_{0})$, che è la
\eqref{eq:convessa}. Per $x=x_{0}$ la disuguaglianza è un'uguaglianza.

\emph{(b)$\Rightarrow$(c).} Se $f'$ è crescente, per ogni $x\in(a,b)$ i rapporti
incrementali $\frac{f'(x+h)-f'(x)}{h}$ sono $\ge0$ (come nella dimostrazione del
\cref{teo:criterio-monotonia}), e il loro limite $f''(x)$ è $\ge0$.

\emph{(c)$\Rightarrow$(b).} Siano $a<c<d<b$. Nell'intervallo $[c,d]$ la funzione $f'$ è
derivabile, quindi continua, e ha derivata $f''\ge0$; per il criterio di monotonia
(\cref{teo:criterio-monotonia}) $f'$ è crescente in $[c,d]$. Poiché due punti qualsiasi
di $(a,b)$ sono contenuti in un tale intervallo, $f'$ è crescente in $(a,b)$. Resta da
includere gli estremi, nei quali $f'$ potrebbe non essere continua. Sia $x\in(a,b)$; per
ogni $h$ con $0<h<x-a$ il teorema di Lagrange in $[a,a+h]$ fornisce un punto
$\xi_{h}\in(a,a+h)$ tale che
\[
  \frac{f(a+h)-f(a)}{h}=f'(\xi_{h})\le f'(x),
\]
dove la disuguaglianza vale perché $\xi_{h}<a+h<x$ e $f'$ è crescente in $(a,b)$. Per
$h\to0^{+}$ il primo membro tende a $f'(a)$, e per il teorema del confronto
$f'(a)\le f'(x)$. Allo stesso modo si prova che $f'(x)\le f'(b)$; infine
$f'(a)\le f'(b)$ si ottiene confrontando entrambi con un qualunque $x\in(a,b)$. Dunque
$f'$ è crescente in $[a,b]$.

Il caso delle funzioni concave si ottiene applicando quanto provato a $-f$.
\end{dimostrazione}

\begin{osservazione}[Convessità stretta]\label{oss:convessita-stretta}
Se $f''(x)>0$ per ogni $x\in(a,b)$, allora $f$ è \emph{strettamente} convessa in $[a,b]$.
Infatti, per il \cref{cor:derivata-positiva} applicato a $f'$, la derivata prima è
strettamente crescente in ogni $[c,d]\subset(a,b)$, e quindi in $(a,b)$; negli estremi,
per $a<t<x$ si ha $f'(a)\le f'(t)<f'(x)$. Nella dimostrazione di (b)$\Rightarrow$(a) le
disuguaglianze tra $f'(x_{1})$ e $f'(x_{0})$ diventano allora strette, e così la
\eqref{eq:convessa} per $x\neq x_{0}$.
\end{osservazione}

Negli enunciati che seguono useremo le derivate successive calcolate in un singolo
punto; precisiamo il significato dell'espressione.

\begin{notazione}[Derivate successive in un punto]\label{not:derivabile-n-volte}
\index{derivata!successiva in un punto}
Si dice che $f$ è \emph{derivabile due volte nel punto $x_{0}$} se $f'$ esiste in tutto un
intorno di $x_{0}$ ed è derivabile in $x_{0}$; la derivata seconda $f''(x_{0})$ è allora la
derivata di $f'$ nel punto $x_{0}$. Più in generale, $f$ è \emph{derivabile $n$ volte in
$x_{0}$} se la derivata $f^{(n-1)}$ esiste in tutto un intorno di $x_{0}$ ed è derivabile in
$x_{0}$. In tal caso $f,f',\dots,f^{(n-2)}$ sono derivabili, e quindi continue, in tale
intorno; nulla si richiede invece a $f^{(n-1)}$ fuori del punto $x_{0}$.
\end{notazione}

\begin{corollario}[Condizione necessaria per i flessi]\label{cor:flesso-necessaria}
Se $x_{0}$ è un punto di flesso per $f$ e $f$ è derivabile due volte in $x_{0}$, allora
$f''(x_{0})=0$.
\end{corollario}

\begin{schema}
\Ipotesi{$x_{0}$ punto di flesso; $f$ derivabile due volte in $x_{0}$.}
\Tesi{$f''(x_{0})=0$.}
\end{schema}

\begin{strategia}
Il cambiamento di concavità si traduce, per il criterio di convessità, in un cambiamento
di monotonia della derivata prima: $f'$ cresce da una parte di $x_{0}$ e decresce
dall'altra. Dunque $x_{0}$ è un punto di estremo per la funzione $f'$, e il teorema di
Fermat, applicato a $f'$, dà $(f')'(x_{0})=0$.
\end{strategia}

\begin{dimostrazione}
Supponiamo $f$ convessa in $[x_{0}-\delta,x_{0}]$ e concava in $[x_{0},x_{0}+\delta]$.
Per il criterio di convessità, nella forma (a)$\Rightarrow$(b) (che non usa la derivata
seconda), $f'$ è crescente in $[x_{0}-\delta,x_{0}]$ e decrescente in
$[x_{0},x_{0}+\delta]$. Dunque $x_{0}$ è un punto di massimo relativo \emph{interno} per
la funzione $f'$, che è derivabile in $x_{0}$: per il teorema di Fermat, applicato a
$f'$, risulta $f''(x_{0})=0$. Il caso opposto è analogo (con un minimo di $f'$).
\end{dimostrazione}

La condizione non è sufficiente: la funzione $f(x)=x^{4}$ ha $f''(0)=0$, ma è convessa su
tutto $\R$, perché $f''(x)=12x^{2}\ge0$; il punto $0$ non è di flesso.

\begin{esempio}[Convessità delle funzioni elementari]\label{es:convessita-elementari}
La funzione $e^{x}$ è convessa su tutto $\R$, perché la sua derivata seconda $e^{x}$ è
positiva; la disuguaglianza $e^{x}\ge1+x$ dell'\cref{es:disuguaglianze-exp-log} è la
\eqref{eq:convessa} con $x_{0}=0$. La funzione $\log x$ è concava in $(0,+\infty)$,
perché $D^{2}\log x=-\frac1{x^{2}}<0$. Per $\operatorname{arctg}x$ si ha
\[
  D\operatorname{arctg}x=\frac1{1+x^{2}}, \qquad
  D^{2}\operatorname{arctg}x=-\frac{2x}{(1+x^{2})^{2}},
\]
quindi l'arcotangente è convessa per $x\le0$ e concava per $x\ge0$, e il punto $0$ è di
flesso (\cref{fig:arctg}). La funzione $x^{2}$ è convessa su $\R$. Infine
$f(x)=x^{3}-3x$ ha $f''(x)=6x$: è concava per $x\le0$ e convessa per $x\ge0$, con un
flesso nell'origine, come si vede nella \cref{fig:cubica-monotonia}.
\end{esempio}

\begin{esempio}[Il grafico di $\sen x$]\label{es:grafico-seno}
Le proprietà stabilite negli ultimi due paragrafi permettono di \emph{motivare} il
grafico, ben noto, della funzione $\sen x$ nell'intervallo $[0,2\pi]$. Le derivate sono
$f'(x)=\cos x$ e $f''(x)=-\sen x$, e i loro segni sono riassunti nella
\cref{tab:seno-segni}. Ne segue che $\sen x$ cresce in $[0,\frac\pi2]$, decresce in
$[\frac\pi2,\frac32\pi]$ e cresce di nuovo in $[\frac32\pi,2\pi]$; è concava in
$[0,\pi]$ e convessa in $[\pi,2\pi]$. Il punto $\frac\pi2$ è di massimo, $\frac32\pi$ di
minimo, $\pi$ di flesso (a tangente obliqua, con pendenza $\cos\pi=-1$). Queste
informazioni, insieme a pochi valori facilmente calcolabili, determinano il grafico
della \cref{fig:seno-studio}. In modo analogo il lettore può studiare $\cos x$.
\end{esempio}

\begin{table}[htb]
\centering
\small
\renewcommand{\arraystretch}{1.35}
\setlength{\tabcolsep}{4.5pt}
\begin{tabular}{l|cccc}
  $f(x)=\sen x$ & $0<x<\frac\pi2$ & $\frac\pi2<x<\pi$ & $\pi<x<\frac32\pi$ &
  $\frac32\pi<x<2\pi$\\\hline
  segno di $f$   & $+$ & $+$ & $-$ & $-$\\
  segno di $f'$  & $+$ & $-$ & $-$ & $+$\\
  segno di $f''$ & $-$ & $-$ & $+$ & $+$\\\hline
  monotonia   & crescente & \multicolumn{2}{c}{decrescente} & crescente\\
  concavità   & \multicolumn{2}{c}{concava} & \multicolumn{2}{c}{convessa}
\end{tabular}
\caption{Segni di $f(x)=\sen x$ e delle sue derivate $f'(x)=\cos x$ e
$f''(x)=-\sen x$ in $(0,2\pi)$, e le conseguenti informazioni di monotonia e di
concavità.}
\label{tab:seno-segni}
\end{table}

\begin{figure}[htb]
\centering
\begin{tikzpicture}
\begin{axis}[width=0.9\textwidth, height=5cm, xmin=-0.3, xmax=7, ymin=-1.35, ymax=1.35,
  xtick={1.5708,3.1416,4.7124,6.2832},
  xticklabels={$\frac\pi2$,$\pi$,{},$2\pi$}, ytick={-1,1},
  xlabel={$x$}, ylabel={$y$}, clip=false]
  \addplot[curva, domain=0:6.2832, samples=150] {sin(deg(x))};
  \addplot[curva2, domain=2.3:4.0, samples=2] {-(x-3.1416)};
  \addplot[tratteggio] coordinates {(1.5708,0) (1.5708,1) (0,1)};
  \addplot[tratteggio] coordinates {(4.7124,0) (4.7124,-1) (0,-1)};
  \addplot[only marks, mark=*, mark size=1.6pt, rame] coordinates
    {(1.5708,1) (3.1416,0) (4.7124,-1)};
  \node[font=\scriptsize, rame, above] at (axis cs:1.5708,1.03) {massimo};
  \node[font=\scriptsize, rame, below] at (axis cs:4.7124,-1.03) {minimo};
  \node[font=\scriptsize, rame, anchor=south west] at (axis cs:3.2,0.08) {flesso};
  \node[font=\scriptsize, anchor=south] at (axis cs:4.7124,0.06) {$\frac32\pi$};
\end{axis}
\end{tikzpicture}
\caption{Il grafico di $\sen x$ in $[0,2\pi]$, ricostruito dalla \cref{tab:seno-segni}.
Nel flesso $x=\pi$ la tangente $y=-(x-\pi)$ attraversa il grafico.}
\label{fig:seno-studio}
\end{figure}

Il programma seguente studia la convessità della funzione $e^{-x^{2}}$, il cui grafico
è la celebre «curva a campana» di Gauss che descrive la distribuzione degli errori
casuali di misura.

\begin{python}[Convessità e flessi della curva di Gauss]
import sympy as sp

x = sp.symbols("x", real=True)
f = sp.exp(-x**2)                       # curva di Gauss
d1 = sp.factor(sp.diff(f, x))
d2 = sp.factor(sp.diff(f, x, 2))
print("f'  =", d1)
print("f'' =", d2)
flessi = sorted(sp.solveset(d2, x, sp.S.Reals))
approx = [round(float(c), 4) for c in flessi]
print("f'' = 0 in", flessi, "~", approx)
for t in (-1, 0, 1):              # segno di f'' negli intervalli
    print(f"  f''({t}) = {d2.subs(x, t)}  ->",
          "convessa" if d2.subs(x, t) > 0 else "concava")
print("f(1/sqrt(2)) =", f.subs(x, flessi[1]),
      "~", round(float(f.subs(x, flessi[1])), 4))
print("punto critico x = 0:  f''(0) =",
      d2.subs(x, 0), "-> massimo relativo")
\end{python}
\begin{risultato}
f'  = -2*x*exp(-x**2)
f'' = 2*(2*x**2 - 1)*exp(-x**2)
f'' = 0 in [-sqrt(2)/2, sqrt(2)/2] ~ [-0.7071, 0.7071]
  f''(-1) = 2*exp(-1)  -> convessa
  f''(0) = -2  -> concava
  f''(1) = 2*exp(-1)  -> convessa
f(1/sqrt(2)) = exp(-1/2) ~ 0.6065
punto critico x = 0:  f''(0) = -2 -> massimo relativo
\end{risultato}

Poiché il fattore $e^{-x^{2}}$ è positivo, il segno di $f''$ è quello di $2x^{2}-1$: la
curva di Gauss è convessa per $\abs x\ge\frac1{\sqrt2}$ e concava per
$\abs x\le\frac1{\sqrt2}$, e i due punti $x=\pm\frac1{\sqrt2}\approx\pm0{,}7071$ sono di
flesso, con ordinata $e^{-1/2}\approx0{,}6065$. L'ultima riga dell'output anticipa il
criterio che dimostriamo alla fine del paragrafo.

\begin{idea}[L'efficienza massima vista con la convessità]
Ritorniamo ancora una volta alla polare parabolica, pensando ora il coefficiente di
resistenza come funzione del coefficiente di portanza:
$C_{D}(C_{L})=C_{D0}+kC_{L}^{2}$. Poiché $C_{D}''=2k>0$, la polare è strettamente
convessa. L'efficienza $E=\frac{C_{L}}{C_{D}}$ è l'inverso della pendenza della retta che
congiunge l'origine con il punto $(C_{L},C_{D}(C_{L}))$; massimizzare $E$ significa
quindi cercare la retta per l'origine di pendenza \emph{minima} che tocca la polare
(\cref{fig:polare-tangente}). La tangente in $C_{L}^{*}$,
$y=C_{D}(C_{L}^{*})+2kC_{L}^{*}(C_{L}-C_{L}^{*})$, passa per l'origine se e solo se
$C_{D0}+kC_{L}^{*2}-2kC_{L}^{*2}=0$, cioè $C_{L}^{*}=\sqrt{C_{D0}/k}$, e in tal caso ha
equazione $y=2kC_{L}^{*}C_{L}=2\sqrt{kC_{D0}}\,C_{L}$. Per la stretta convessità la polare
sta sopra questa tangente:
\[
  C_{D}(C_{L})>2\sqrt{kC_{D0}}\;C_{L} \quad\text{per } C_{L}\neq C_{L}^{*},
\]
e quindi, dividendo per $C_{L}\,C_{D}(C_{L})>0$,
\[
  E(C_{L})=\frac{C_{L}}{C_{D}(C_{L})}<\frac1{2\sqrt{kC_{D0}}}=E_{\max}
  \qquad\text{per } C_{L}>0,\ C_{L}\neq C_{L}^{*}.
\]
Ritroviamo il risultato del \cref{cap:funzioni-continue} con una dimostrazione puramente
geometrica, che è anche la costruzione grafica usata dagli aerodinamici: \emph{l'assetto
di massima efficienza si trova tracciando dall'origine la tangente alla polare}.
\end{idea}

\begin{figure}[htb]
\centering
\begin{tikzpicture}
\begin{axis}[width=0.85\textwidth, height=6cm, xmin=0, xmax=1.62, ymin=0, ymax=0.15,
  axis lines=left, xlabel={$C_{L}$}, ylabel={$C_{D}$}, xtick={0.709,1.2},
  xticklabels={$C_{L}^{*}$,$C_{L}$}, ytick={0.025,0.05},
  yticklabels={$C_{D0}$,$0{,}05$}, clip=false,
  scaled y ticks=false]
  \addplot[curva, domain=0:1.55, samples=100] {0.025+0.0497359*x^2};
  \addplot[curva2, domain=0:1.55, samples=2] {x/14.17963};
  \addplot[curva3, domain=0:1.4, samples=2] {0.0805164*x};
  \addplot[tratteggio] coordinates {(0.709,0) (0.709,0.05)};
  \addplot[tratteggio] coordinates {(1.2,0) (1.2,0.09662)};
  \addplot[only marks, mark=*, mark size=1.6pt, rame] coordinates {(0.709,0.05)};
  \addplot[only marks, mark=*, mark size=1.6pt, verde] coordinates {(1.2,0.09662)};
  \node[font=\scriptsize, blu, anchor=east] at (axis cs:1.38,0.13) {polare};
  \node[font=\scriptsize, rame, anchor=west] at (axis cs:1.56,0.108)
    {pendenza $\frac1{E_{\max}}$};
  \node[font=\scriptsize, verde!50!black, anchor=south east] at (axis cs:1.15,0.098)
    {pendenza $\frac1{E(C_{L})}$};
\end{axis}
\end{tikzpicture}
\caption{La polare parabolica $C_{D}=C_{D0}+kC_{L}^{2}$ per $C_{D0}=0{,}025$, $AR=8$,
$e=0{,}8$. Ogni retta per l'origine che incontra la polare ha pendenza $\frac1E$; la
tangente dall'origine (rame) ha la pendenza minima e tocca la polare in
$C_{L}^{*}\approx0{,}709$, dove $E=E_{\max}\approx14{,}18$. La secante verde, per
$C_{L}=1{,}2$, ha pendenza maggiore: $E\approx12{,}4$.}
\label{fig:polare-tangente}
\end{figure}

\subsection{Il criterio della derivata seconda}

Chiudiamo il paragrafo con un criterio, basato sul segno della derivata seconda, per
stabilire se un punto stazionario è di massimo o di minimo relativo. Lo generalizzeremo
nel \cref{sec:taylor}.

\begin{teorema}[Criterio della derivata seconda]\label{teo:test-derivata-seconda}
\index{criterio!della derivata seconda}
Sia $f$ derivabile due volte in un intorno di $x_{0}$, con derivata seconda continua in
$x_{0}$. Allora
\begin{equation}\label{eq:test-derivata-seconda}
  \begin{aligned}
  &f'(x_{0})=0,\ f''(x_{0})>0 \ \implica\ x_{0} \text{ di minimo relativo stretto};\\
  &f'(x_{0})=0,\ f''(x_{0})<0 \ \implica\ x_{0} \text{ di massimo relativo stretto}.
  \end{aligned}
\end{equation}
\end{teorema}

\begin{schema}
\Ipotesi{$f$ derivabile due volte vicino a $x_{0}$, $f''$ continua in $x_{0}$;
$f'(x_{0})=0$ e $f''(x_{0})>0$ (risp.\ $<0$).}
\Tesi{$x_{0}$ è di minimo (risp.\ massimo) relativo stretto.}
\end{schema}

\begin{strategia}
Se $f''(x_{0})>0$ e $f''$ è continua, per la permanenza del segno $f''$ è positiva in
tutto un intorno di $x_{0}$; lì la funzione è strettamente convessa e il grafico sta
sopra la tangente in $x_{0}$. Ma la tangente in un punto stazionario è la retta
orizzontale $y=f(x_{0})$: dunque il grafico sta sopra il livello $f(x_{0})$, che è la
definizione di minimo relativo.
\end{strategia}

\begin{dimostrazione}
Sia $f''(x_{0})>0$. Per il teorema della permanenza del segno per le funzioni continue
(\cref{teo:permanenza-segno-continue}) esiste $\delta>0$ tale che $f''(x)>0$ per ogni
$x\in(x_{0}-\delta,x_{0}+\delta)$, intervallo che possiamo supporre contenuto
nell'intorno in cui $f$ è derivabile due volte. Fissato $0<\delta'<\delta$, per il
criterio di convessità e l'\cref{oss:convessita-stretta}, $f$ è strettamente convessa in
$[x_{0}-\delta',x_{0}+\delta']$, e quindi, ricordando che $f'(x_{0})=0$,
\begin{equation}\label{eq:test-derivata-seconda-dim}
  f(x)>f(x_{0})+f'(x_{0})(x-x_{0})=f(x_{0})
\end{equation}
per ogni $x\in[x_{0}-\delta',x_{0}+\delta']$ con $x\neq x_{0}$.
Dunque $x_{0}$ è un punto di minimo relativo stretto. Il caso $f''(x_{0})<0$ si ottiene
applicando quanto provato a $-f$.
\end{dimostrazione}

\begin{osservazione}[Basta l'esistenza di $f''(x_{0})$]\label{oss:test-derivata-seconda}
La continuità di $f''$ in $x_{0}$ non è necessaria. Supponiamo soltanto che $f$ sia
derivabile in un intorno di $x_{0}$, con $f'(x_{0})=0$, e che esista $f''(x_{0})>0$.
Allora
\[
  f''(x_{0})=\lim_{x\to x_{0}}\frac{f'(x)-f'(x_{0})}{x-x_{0}}=\lim_{x\to x_{0}}\frac{f'(x)}{x-x_{0}}>0,
\]
e per la permanenza del segno (\cref{teo:confronto-funzioni}\,(i)) il rapporto
$\frac{f'(x)}{x-x_{0}}$ è positivo per $0<\abs{x-x_{0}}<\delta$. Quindi $f'(x)<0$ per
$x_{0}-\delta<x<x_{0}$ e $f'(x)>0$ per $x_{0}<x<x_{0}+\delta$, e il test della derivata
prima (\cref{cor:test-derivata-prima}) mostra che $x_{0}$ è di minimo relativo stretto.
Per la curva di Gauss, $f'(0)=0$ e $f''(0)=-2<0$: l'origine è di massimo relativo
(anzi assoluto, perché $e^{-x^{2}}\le1$). Se invece $f'(x_{0})=f''(x_{0})=0$, il criterio
non decide: $x^{4}$, $-x^{4}$ e $x^{3}$ hanno in $0$, rispettivamente, un minimo, un
massimo e un flesso. Il \cref{sec:taylor} risolverà questi casi.
\end{osservazione}

\subsection{Approfondimento: la convessità senza derivate}
\label{subsec:convessita-corde}

La \cref{def:convessa} confronta il grafico con le tangenti, e richiede quindi la
derivabilità; eppure anche $\abs x$, che ha uno spigolo nell'origine, «volge la concavità
verso l'alto». Seguendo \cite{giusti}, usiamo le \emph{corde} al posto delle tangenti:
per le funzioni derivabili si ritrova la nostra definizione, e dalla sola convessità
discendono la continuità e l'esistenza delle derivate unilaterali nei punti interni.

\begin{definizione}[Convessità nel senso delle corde]\label{def:convessa-corde}
\index{funzione!convessa!nel senso delle corde}
Sia $I$ un intervallo e sia $f:I\to\R$. Si dice che $f$ è \emph{convessa nel senso delle
corde} in $I$ se, per ogni $x_{1},x_{2}\in I$ e per ogni $t\in[0,1]$, risulta
\begin{equation}\label{eq:convessa-corde}
  f\bigl((1-t)x_{1}+tx_{2}\bigr)\le(1-t)f(x_{1})+tf(x_{2}) ;
\end{equation}
si dice che $f$ è \emph{concava nel senso delle corde} se $-f$ è convessa nel senso
delle corde.
\end{definizione}

Se $x_{1}<x_{2}$, al variare di $t$ in $[0,1]$ il punto $x=(1-t)x_{1}+tx_{2}$ descrive
$[x_{1},x_{2}]$, con $t=\frac{x-x_{1}}{x_{2}-x_{1}}$, e la condizione diventa
\begin{equation}\label{eq:corda-sopra}
  f(x)\le f(x_{1})+\frac{f(x_{2})-f(x_{1})}{x_{2}-x_{1}}\,(x-x_{1})
  \qquad (x_{1}\le x\le x_{2}) :
\end{equation}
sopra $[x_{1},x_{2}]$ il grafico sta sotto la corda di estremi $(x_{1},f(x_{1}))$ e
$(x_{2},f(x_{2}))$. Il caso $x_{1}=x_{2}$ è banale, e scambiare $x_{1}$ con $x_{2}$
equivale a sostituire $t$ con $1-t$. Poiché la condizione coinvolge solo i punti di
$[x_{1},x_{2}]$, $f$ è convessa nel senso delle corde in $I$ se e solo se lo è in ogni
intervallo chiuso e limitato contenuto in $I$. Ad esempio $\abs x$ lo è su $\R$, perché
$\abs{(1-t)x_{1}+tx_{2}}\le(1-t)\abs{x_{1}}+t\abs{x_{2}}$ per la disuguaglianza
triangolare (\cref{teo:triangolare}).

\begin{lemma}[delle tre corde]\label{lem:tre-corde}\index{lemma delle tre corde}
Sia $f$ convessa nel senso delle corde nell'intervallo $I$.
\begin{enumerate}[label=\textup{(\roman*)}]
  \item Se $x_{1}<x<x_{2}$ sono punti di $I$, allora (\cref{fig:tre-corde})
  \begin{equation}\label{eq:tre-corde}
    \frac{f(x)-f(x_{1})}{x-x_{1}}\le\frac{f(x_{2})-f(x_{1})}{x_{2}-x_{1}}
    \le\frac{f(x_{2})-f(x)}{x_{2}-x} .
  \end{equation}
  \item Per ogni $x_{0}\in I$ il rapporto incrementale
  $R(x)=\frac{f(x)-f(x_{0})}{x-x_{0}}$ è una funzione crescente in $I\setminus\{x_{0}\}$.
\end{enumerate}
\end{lemma}

\begin{schema}
\Ipotesi{$f$ convessa nel senso delle corde nell'intervallo $I$.}
\Tesi{(i) per $x_{1}<x<x_{2}$ in $I$ vale la \eqref{eq:tre-corde}; \ (ii) per ogni
$x_{0}\in I$, $R(x)=\frac{f(x)-f(x_{0})}{x-x_{0}}$ è crescente in $I\setminus\{x_{0}\}$.}
\end{schema}

\begin{strategia}
Il punto $P=(x,f(x))$ sta sotto la retta $P_{1}P_{2}$: abbassando l'estremo destro di un
segmento che parte da $P_{1}$ la pendenza diminuisce, abbassando l'estremo sinistro di un
segmento che arriva in $P_{2}$ la pendenza aumenta. Per la (ii) si applica la (i) alla
terna formata da $x_{0}$ e dai due punti da confrontare.
\end{strategia}

\begin{dimostrazione}
(i) Sia $m=\frac{f(x_{2})-f(x_{1})}{x_{2}-x_{1}}$ e sia $r(s)=f(x_{1})+m(s-x_{1})$, cosicché
$r(x_{1})=f(x_{1})$, $r(x_{2})=f(x_{2})$ e, per la \eqref{eq:corda-sopra}, $f(x)\le r(x)$.
Poiché $x-x_{1}>0$ e $x_{2}-x>0$,
\[
\begin{gathered}
  \frac{f(x)-f(x_{1})}{x-x_{1}}\le\frac{r(x)-r(x_{1})}{x-x_{1}}=m ,\\[2pt]
  \frac{f(x_{2})-f(x)}{x_{2}-x}\ge\frac{r(x_{2})-r(x)}{x_{2}-x}=m ,
\end{gathered}
\]
perché $r$ è un polinomio di primo grado di pendenza $m$: è la \eqref{eq:tre-corde}.

(ii) Il rapporto incrementale non cambia scambiando i due punti:
$R(x)=\frac{f(x_{0})-f(x)}{x_{0}-x}$. Siano $x<y$ due punti di $I\setminus\{x_{0}\}$. Se
$x_{0}<x<y$, la prima disuguaglianza della \eqref{eq:tre-corde}, scritta per la terna
$x_{0}<x<y$, dice che $R(x)\le R(y)$; se $x<x_{0}<y$, lo dice il confronto tra il primo e
il terzo membro per la terna $x<x_{0}<y$; se $x<y<x_{0}$, lo dice la seconda
disuguaglianza per la terna $x<y<x_{0}$. Dunque $R$ è crescente.
\end{dimostrazione}

\begin{figure}[htb]
\centering
\begin{tikzpicture}
\begin{axis}[width=0.8\textwidth, height=5cm, xmin=0, xmax=4.1, ymin=0, ymax=4.9,
  axis lines=left, xtick={0.5,2,3.5}, xticklabels={$x_{1}$,$x$,$x_{2}$},
  ytick=\empty, clip=false]
  \addplot[curva, domain=0.2:3.7, samples=120] {0.35*x^2+0.3};
  \addplot[tratteggio] coordinates {(0.5,0) (0.5,0.3875)};
  \addplot[tratteggio] coordinates {(2,0) (2,1.7)};
  \addplot[tratteggio] coordinates {(3.5,0) (3.5,4.5875)};
  \draw[curva2] (axis cs:0.5,0.3875) -- (axis cs:2,1.7);
  \draw[curva3] (axis cs:0.5,0.3875) -- (axis cs:3.5,4.5875);
  \draw[grigio, line width=1pt] (axis cs:2,1.7) -- (axis cs:3.5,4.5875);
  \addplot[only marks, mark=*, mark size=1.6pt, blu] coordinates
    {(0.5,0.3875) (2,1.7) (3.5,4.5875)};
  \node[font=\footnotesize, rame, anchor=south] at (axis cs:1.75,1.56) {$s_{1}$};
  \node[font=\footnotesize, verde!50!black, anchor=south east] at (axis cs:1.45,1.85)
    {$s$};
  \node[font=\footnotesize, grigio, anchor=south] at (axis cs:2.2,2.2) {$s_{2}$};
  \node[font=\scriptsize, anchor=south east] at (axis cs:0.47,0.42) {$P_{1}$};
  \node[font=\scriptsize, anchor=north west] at (axis cs:2.04,1.66) {$P$};
  \node[font=\scriptsize, anchor=north west] at (axis cs:3.54,4.55) {$P_{2}$};
\end{axis}
\end{tikzpicture}
\caption{Il lemma delle tre corde. Se $f$ è convessa nel senso delle corde e
$x_{1}<x<x_{2}$, il punto $P$ sta sotto la corda $P_{1}P_{2}$, e le pendenze delle corde
$P_{1}P$, $P_{1}P_{2}$, $PP_{2}$ crescono: $s_{1}\le s\le s_{2}$. Qui
$f(x)=0{,}35\,x^{2}+0{,}3$, con $s_{1}=0{,}875$, $s=1{,}4$ e $s_{2}=1{,}925$.}
\label{fig:tre-corde}
\end{figure}

\begin{proposizione}[Le due definizioni coincidono per le funzioni derivabili]
\label{prop:convessita-equivalenti}
Sia $f$ derivabile in $[a,b]$ (con derivate unilaterali negli estremi). Allora $f$ è
convessa in $[a,b]$ secondo la \cref{def:convessa} se e solo se è convessa nel senso
delle corde in $[a,b]$.
\end{proposizione}

\begin{schema}
\Ipotesi{$f$ derivabile in $[a,b]$.}
\Tesi{grafico sopra le tangenti, \eqref{eq:convessa} \ $\sse$ \ grafico sotto le corde,
\eqref{eq:convessa-corde}.}
\end{schema}

\begin{strategia}
Dalle tangenti alle corde: si scrive la \eqref{eq:convessa} con il punto di tangenza nel
punto $z$ in cui si stima $f$, e si combinano le disuguaglianze relative agli estremi
della corda. Dalle corde alle tangenti: il rapporto incrementale $R$, crescente per il
lemma delle tre corde, a destra di $x_{0}$ non è minore del suo limite $f'(x_{0})$, a
sinistra non è maggiore; si moltiplica poi per $x-x_{0}$.
\end{strategia}

\begin{dimostrazione}
($\Rightarrow$) È l'argomento dell'\cref{es:corde-convesse}. Siano $x_{1},x_{2}\in[a,b]$,
$t\in[0,1]$ e $z=(1-t)x_{1}+tx_{2}$. Per la \eqref{eq:convessa} con punto di tangenza $z$,
\[
  f(x_{1})\ge f(z)+f'(z)(x_{1}-z), \qquad f(x_{2})\ge f(z)+f'(z)(x_{2}-z) .
\]
Moltiplicando la prima per $1-t\ge0$, la seconda per $t\ge0$ e sommando, si ottiene
$(1-t)f(x_{1})+tf(x_{2})\ge f(z)+f'(z)\bigl[(1-t)x_{1}+tx_{2}-z\bigr]=f(z)$.

($\Leftarrow$) Siano $x_{0},x\in[a,b]$, $x\neq x_{0}$, e sia $R$ il rapporto incrementale
del \cref{lem:tre-corde}\,(ii). Se $x>x_{0}$, per ogni $s\in(x_{0},x)$ si ha
$R(s)\le R(x)$; per $s\to x_{0}^{+}$ il primo membro tende a $f'(x_{0})$
(\eqref{eq:derivata-xi} e \cref{prop:dx-sx}; se $x_{0}=a$ si tratta della derivata
destra), e per il teorema del confronto (\cref{teo:confronto-funzioni}\,(ii)) risulta
$f'(x_{0})\le R(x)$. Moltiplicando per $x-x_{0}>0$ si ottiene
$f(x)-f(x_{0})\ge f'(x_{0})(x-x_{0})$. Se $x<x_{0}$, allo stesso modo $R(x)\le f'(x_{0})$,
e moltiplicando per $x-x_{0}<0$ si ritrova la stessa disuguaglianza, cioè la
\eqref{eq:convessa}.
\end{dimostrazione}

Dunque $x^{2}$, $e^{x}$, $-\log x$ e tutte le funzioni convesse incontrate finora lo sono
anche nel senso delle corde. Una funzione convessa nel senso delle corde può non essere
derivabile, come $\abs x$; ma nei punti interni non può essere troppo irregolare.

\begin{teorema}[Regolarità delle funzioni convesse]\label{teo:convesse-continue}
Sia $f$ convessa nel senso delle corde nell'intervallo $I$, e sia $x_{0}$ un punto
interno di $I$. Allora:
\begin{enumerate}[label=\textup{(\roman*)}]
  \item esistono finite la derivata sinistra $f'_{-}(x_{0})$ e la derivata destra
  $f'_{+}(x_{0})$, e risulta $f'_{-}(x_{0})\le f'_{+}(x_{0})$;
  \item $f$ è continua in $x_{0}$;
  \item per ogni $m\in\bigl[f'_{-}(x_{0}),f'_{+}(x_{0})\bigr]$ si ha
  \begin{equation}\label{eq:retta-appoggio}
    f(x)\ge f(x_{0})+m(x-x_{0}) \qquad\perogni x\in I .
  \end{equation}
\end{enumerate}
\end{teorema}

\begin{schema}
\Ipotesi{$f$ convessa nel senso delle corde nell'intervallo $I$; \ $x_{0}$ interno a $I$.}
\Tesi{(i) $f'_{-}(x_{0})\le f'_{+}(x_{0})$, entrambe finite; \ (ii) $f$ continua in
$x_{0}$; \ (iii) vale la \eqref{eq:retta-appoggio}.}
\end{schema}

\begin{strategia}
Tutto discende dalla monotonia di $R$. Una funzione monotòna ha limiti unilaterali, finiti
se è limitata dalla parte giusta: qui interviene l'ipotesi che $x_{0}$ sia interno, perché
un punto $a'<x_{0}$ fornisce il minorante $R(a')$ a destra di $x_{0}$, e un punto
$b'>x_{0}$ un maggiorante a sinistra. La continuità si legge nell'identità
$f(x)-f(x_{0})=R(x)(x-x_{0})$; la (iii) si ottiene confrontando $R$ con i suoi limiti.
\end{strategia}

\begin{dimostrazione}
Siano $a',b'\in I$ con $a'<x_{0}<b'$; per il \cref{lem:tre-corde}, $R$ è crescente in
$I\setminus\{x_{0}\}$.

(i) La funzione che vale $R(a')$ in $x_{0}$ e coincide con $R$ in $(x_{0},b']$ è crescente
in $[x_{0},b']$, perché $R(a')\le R(x)$ per $x>x_{0}$. Per il \cref{teo:limite-monotone-funzioni}
esiste finito il limite di $R(x)$ per $x\to x_{0}^{+}$, che per la \eqref{eq:derivata-xi}
è $f'_{+}(x_{0})$. Allo stesso modo, la funzione che vale $R(b')$ in $x_{0}$ e coincide
con $R$ in $[a',x_{0})$ è crescente in $[a',x_{0}]$, e quindi esiste finita
$f'_{-}(x_{0})=\lim_{x\to x_{0}^{-}}R(x)$. Infine, se $x<x_{0}<y$, da $R(x)\le R(y)$ si
ottiene, per $x\to x_{0}^{-}$, $f'_{-}(x_{0})\le R(y)$ e poi, per $y\to x_{0}^{+}$,
$f'_{-}(x_{0})\le f'_{+}(x_{0})$ (\cref{teo:confronto-funzioni}\,(ii)).

(ii) Per $x\to x_{0}^{\pm}$ il fattore $R(x)$ ha limite finito e $x-x_{0}\to0$, quindi
$f(x)-f(x_{0})=R(x)(x-x_{0})\to0$ da entrambe le parti; per la \cref{prop:dx-sx}, $f$ è
continua in $x_{0}$.

(iii) Se $x>x_{0}$, per ogni $s\in(x_{0},x)$ è $R(s)\le R(x)$, e al limite per
$s\to x_{0}^{+}$ si ha $m\le f'_{+}(x_{0})\le R(x)$; moltiplicando per $x-x_{0}>0$ si
ottiene la \eqref{eq:retta-appoggio}. Se $x<x_{0}$, allo stesso modo
$R(x)\le f'_{-}(x_{0})\le m$, e si moltiplica per $x-x_{0}<0$. Per $x=x_{0}$ vale
l'uguaglianza.
\end{dimostrazione}

\begin{osservazione}[Rette d'appoggio, spigoli ed estremi dell'intervallo]
\label{oss:convesse-esempi}\index{retta d'appoggio}
La retta $y=f(x_{0})+m(x-x_{0})$ della \eqref{eq:retta-appoggio} si dice \emph{retta
d'appoggio}: passa per $(x_{0},f(x_{0}))$ e lascia tutto il grafico da una parte. Se $f$ è
derivabile in $x_{0}$ l'unica retta d'appoggio è la tangente; in uno spigolo le rette
d'appoggio sono infinite: per $\abs x$ si ha $f'_{-}(0)=-1$, $f'_{+}(0)=1$, e
$\abs x\ge mx$ per ogni $m\in[-1,1]$. Se $f$ è convessa e derivabile in un punto interno
$x_{0}$ con $f'(x_{0})=0$, la \eqref{eq:retta-appoggio} con $m=0$ dà $f(x)\ge f(x_{0})$
per ogni $x\in I$: \emph{per una funzione convessa ogni punto stazionario interno è di
minimo assoluto}. Infine, l'ipotesi che $x_{0}$ sia interno non si può togliere: la
funzione che vale $0$ in $[0,1)$ e $1$ in $1$ è convessa nel senso delle corde in $[0,1]$
(le corde hanno ordinate $\ge0$), ma è discontinua nell'estremo $1$.
\end{osservazione}

La \eqref{eq:convessa-corde} riguarda medie pesate di \emph{due} punti; per induzione la
si estende a un numero qualunque di punti.

\begin{teorema}[Disuguaglianza di Jensen]\label{teo:jensen}\index{disuguaglianza di Jensen}
Sia $f$ convessa nel senso delle corde nell'intervallo $I$, siano $x_{1},\dots,x_{n}\in I$
e siano $t_{1},\dots,t_{n}\ge0$ numeri con $t_{1}+\dots+t_{n}=1$. Allora il punto
$\sum_{k=1}^{n}t_{k}x_{k}$ appartiene a $I$ e
\begin{equation}\label{eq:jensen}
  f\Bigl(\sum_{k=1}^{n}t_{k}x_{k}\Bigr)\le\sum_{k=1}^{n}t_{k}f(x_{k}) .
\end{equation}
\end{teorema}

\begin{schema}
\Ipotesi{$f$ convessa nel senso delle corde in $I$; \ $x_{k}\in I$, $t_{k}\ge0$,
$\sum_{k}t_{k}=1$.}
\Tesi{$\sum_{k}t_{k}x_{k}\in I$ \ e \
$f\bigl(\sum_{k}t_{k}x_{k}\bigr)\le\sum_{k}t_{k}f(x_{k})$.}
\end{schema}

\begin{strategia}
Per induzione su $n$: per $n=2$ la tesi è la definizione. Per passare da $n$ a $n+1$ punti
si raggruppano i primi $n$ nella loro media pesata $y$, con i pesi divisi per la loro
somma: la combinazione diventa una combinazione dei \emph{due} punti $y$ e $x_{n+1}$, a cui
si applica la definizione, mentre $f(y)$ si stima con l'ipotesi induttiva. In termini
meccanici: il baricentro di masse poste sul grafico sta sopra il grafico.
\end{strategia}

\begin{dimostrazione}
Se $p$ e $q$ sono il minimo e il massimo dei punti $x_{k}$, risulta
$p=\sum_{k}t_{k}p\le\sum_{k}t_{k}x_{k}\le\sum_{k}t_{k}q=q$; poiché $I$ è un intervallo che
contiene $p$ e $q$, contiene anche $\sum_{k}t_{k}x_{k}$.

Proviamo la \eqref{eq:jensen} per induzione (\cref{teo:induzione}). Per $n=1$ è $t_{1}=1$
e i due membri coincidono. Supponiamola vera per $n$ punti e siano dati $n+1$ punti con
pesi $t_{1},\dots,t_{n+1}$; sia $s=t_{1}+\dots+t_{n}=1-t_{n+1}$. Se $s=0$ i due membri
valgono $f(x_{n+1})$. Se $s>0$, il punto $y=\sum_{k=1}^{n}\frac{t_{k}}{s}x_{k}$ è una media
pesata con pesi $\ge0$ di somma $1$: quindi $y\in I$ e, per l'ipotesi induttiva,
$f(y)\le\sum_{k=1}^{n}\frac{t_{k}}{s}f(x_{k})$. Poiché
$\sum_{k=1}^{n+1}t_{k}x_{k}=s\,y+t_{n+1}x_{n+1}$, la \eqref{eq:convessa-corde} (con
$t=t_{n+1}$) dà
\[
\begin{aligned}
  f\Bigl(\sum_{k=1}^{n+1}t_{k}x_{k}\Bigr)&\le s\,f(y)+t_{n+1}f(x_{n+1})\\
  &\le\sum_{k=1}^{n}t_{k}f(x_{k})+t_{n+1}f(x_{n+1})=\sum_{k=1}^{n+1}t_{k}f(x_{k}) .
\end{aligned}
\]
La tesi vale dunque per $n+1$ punti, e quindi per ogni $n$.
\end{dimostrazione}

\begin{corollario}[Media aritmetica e media geometrica]
\label{cor:media-aritmetica-geometrica}
\index{disuguaglianza tra media aritmetica e geometrica}
Per ogni $n\in\N$ e per ogni scelta di numeri positivi $a_{1},\dots,a_{n}$,
\begin{equation}\label{eq:am-gm}
  \sqrt[n]{a_{1}a_{2}\cdots a_{n}}\le\frac{a_{1}+a_{2}+\dots+a_{n}}{n} .
\end{equation}
\end{corollario}

\begin{schema}
\Ipotesi{$n\in\N$; \ $a_{1},\dots,a_{n}>0$.}
\Tesi{la media geometrica non supera la media aritmetica: vale la \eqref{eq:am-gm}.}
\end{schema}

\begin{strategia}
Il logaritmo della media geometrica è la media aritmetica dei logaritmi. La tesi diventa
così «il logaritmo della media non è minore della media dei logaritmi»: è la
disuguaglianza di Jensen per la funzione convessa $-\log x$, con i pesi $t_{k}=\frac1n$.
\end{strategia}

\begin{dimostrazione}
La funzione $\varphi(x)=-\log x$ ha $\varphi''(x)=\frac1{x^{2}}>0$ in $(0,+\infty)$
(\cref{es:convessita-elementari}); per il \cref{teo:criterio-convessita} e la
\cref{prop:convessita-equivalenti} è convessa nel senso delle corde in ogni intervallo
$[a,b]\subset(0,+\infty)$, e quindi in $(0,+\infty)$. La \eqref{eq:jensen} con
$x_{k}=a_{k}$ e $t_{k}=\frac1n$ e le proprietà dei logaritmi (\cref{prop:log-proprieta})
danno
\[
  -\log\Bigl(\frac1n\sum_{k=1}^{n}a_{k}\Bigr)\le-\frac1n\sum_{k=1}^{n}\log a_{k}
  =-\log\sqrt[n]{a_{1}a_{2}\cdots a_{n}} .
\]
Cambiando segno, il logaritmo della media geometrica non supera quello della media
aritmetica; poiché il logaritmo è strettamente crescente (\cref{prop:log-monotono}), lo
stesso vale per le due medie.
\end{dimostrazione}

Per $n=2$ si ritrova la disuguaglianza $\sqrt{ab}\le\frac{a+b}2$, dimostrata per via
algebrica negli esercizi del \cref{cap:numeri}. Se gli $a_{k}$ sono tutti uguali, i due
membri della \eqref{eq:am-gm} coincidono; con la convessità stretta di $-\log x$ si
dimostra che è l'unico caso di uguaglianza.

\begin{esempio}[La potenza minima senza derivate]\label{es:potenza-am-gm}
Con termini scelti ad arte, la \eqref{eq:am-gm} trova minimi senza derivate: si maggiora
(o si minora) la funzione con una costante, e si cerca dove la disuguaglianza diventa
un'uguaglianza. Nel \cref{sec:studio-funzione} vedremo che la potenza necessaria al volo
livellato è $P(V)=aV^{3}+\frac bV$, con $a,b>0$. Spezzando $\frac bV$ in tre parti uguali,
in modo che nel prodotto dei quattro termini $V$ si semplifichi,
\[
  P(V)=aV^{3}+3\cdot\frac{b}{3V}\ge4\sqrt[4]{aV^{3}\Bigl(\frac{b}{3V}\Bigr)^{3}}
  =4\sqrt[4]{\frac{ab^{3}}{27}} ,
\]
con l'uguaglianza quando i quattro termini sono uguali: $aV^{3}=\frac{b}{3V}$, cioè
$V^{4}=\frac{b}{3a}$, la velocità di minima potenza \eqref{eq:V-min-P}. Per il velivolo
del \cref{sec:studio-funzione} si trova $P_{\min}\approx23{,}47$~kW, come là con le
derivate. Con due soli termini lo stesso procedimento dà la resistenza minima e, nel
\cref{cap:funzioni-continue}, l'efficienza massima.
\end{esempio}

\begin{notastorica}[Le funzioni convesse]
Le funzioni convesse comparvero in modo implicito in molti lavori dell'Ottocento sulle
disuguaglianze tra medie, ma la loro teoria sistematica fu iniziata dal matematico
danese Johan Jensen, che nel 1906 pubblicò negli \emph{Acta Mathematica} la memoria
\emph{Sur les fonctions convexes et les inégalités entre les valeurs moyennes}. Jensen
adottò una definizione che non richiede la derivabilità, basata sui punti medi delle
corde anziché sulle tangenti: $f\bigl(\frac{x+y}2\bigr)\le\frac{f(x)+f(y)}2$ per ogni $x,y$.
Per le funzioni continue essa equivale alla \eqref{eq:convessa-corde}, cioè al fatto che
ogni corda stia sopra il grafico, che è la definizione oggi più comune: per le funzioni
derivabili coincide con la nostra (\cref{prop:convessita-equivalenti}), ma si applica
anche a funzioni non derivabili, come $\abs x$. Nella stessa memoria Jensen dimostrò in
generale la disuguaglianza \eqref{eq:jensen}, già ottenuta da Otto Hölder nel 1889 per
le funzioni con $f''\ge0$. Oggi la convessità è al centro dell'\emph{ottimizzazione
convessa}, la teoria che garantisce che un minimo locale sia anche globale e che sta
alla base di molti algoritmi di progetto ottimo in ingegneria.
\end{notastorica}
